\documentclass[12 pt]{article}
\usepackage[a4paper,
 mag=1000, includefoot,
 left=3cm, right=1.5cm, top=2cm, bottom=2cm, headsep=1cm, footskip=1cm]{geometry}
\usepackage[T2A]{fontenc}
\usepackage[utf8]{inputenc}
\usepackage[english,russian]{babel}
\usepackage{indentfirst}
\usepackage{amsmath}
\usepackage{array}
\usepackage{graphicx}
\usepackage{hyperref}
\usepackage{footnote}
\usepackage{breakurl}
\usepackage{url}
\usepackage{amssymb}
\usepackage{xcolor}
\usepackage{tipa}
\usepackage{upgreek}
\usepackage{listings}
\usepackage{longtable}
\usepackage{caption}
\usepackage{textcomp}
\usepackage{marvosym}
\usepackage{amsthm}
\newtheorem{mydef}{Определение}
\newtheorem{myth}{Теорема}
\newtheorem{mynote}{Замечание}
\begin{document}

Гильбертово пространство --- (полное) линейное пространство со скалярным произведением. Например, это обычное евклидово пространство векторов, в котором обычное скалярное произведение $\langle X,Y \rangle = \sum_i x_i y_i$. Также, можно задать положительные веса $w_i$ и рассмотреть взвешенное скалярное произведение $\langle X,Y \rangle = \sum_i w_i x_i y_i$. Получится другое гильбертово пространство из векторов с весами.

Посмотрим на векторы с другой стороны. Пусть $M=\{1,\ldots,L\}$ – множество чисел от $1$ до $L$. Тогда вектор $X$ можно рассматривать как отображение, заданное на $M$, со значениями в $\mathbb{R}$, так как компоненту вектора $x_i$ можно рассматривать как результат применения отображения $X$ к $i$: $X(i)=x_i$.

Посмотрим на евклидово скалярное произведение с другой стороны.
Для этого нам понадобится понятие интеграла по мере. Зададим меру $\mu$ на $M$: $\mu(\{i\}) = w_i>0$.
Тогда интегралом по дискретной мере от функции является сумма $\int_M f(i) \mu(d\,i)= \sum f(i) w_i$.
Вообще, интеграл по мере $\int_D f(x) \mu(d\,x)$ определяется в общем случае. Если мера дискретная (как выше), то
интеграл превращается в сумму. Если $D$ --- все $\mathbb{R}$ (или его интервал), а $\mu(dx) = w(x) dx$, где $w(x)>0$,
то $\int_D f(x) \mu(d\,x) = \int_D f(x) w(x) d\,x$, т.е. обычный интеграл.

Через введенное понятие интеграла по мере можем записать
$$\langle X,Y \rangle = \int_M X(i) Y(i) \mu(d\,i) = \sum_i w_i x_i y_i.$$

Теперь построим следующую конструкцию. Пусть $(D,\mu)$ --- измеримое пространство (там еще должна быть
сигма-алгебра, где задана мера, но мы ее опустим).
Рассмотрим гильбертово пространство функций $L^2(D,\mu)$, заданных на $(D,\mu)$, со скалярным произведением
$$\langle f_1,f_2 \rangle = \int_D f_1(x) f_2(x) \mu(d\,x).$$
В это пространство входят не все функции, а только те, чьи нормы конечны, т.е.
$||f||^2 = \langle f,f\rangle <\infty$.

\textit{Задание 1: какое должно быть пространство $D$ и какая должна быть мера $\mu$, чтобы $L^2(D,\mu)$
было обычным евклидовым пространством из векторов с обычным скалярным произведением?}

Итак, мы умеем задавать вектора как функции и их скалярное произведение как интеграл по мере.

Теперь пусть у нас есть два гильбертовых пространства функций, заданных на $(D_1,\mu_1)$ и
$(D_2,\mu_2)$. Обозначим их $L_1 = L^2(D_1,\mu_1)$ и  $L_2 = L^2(D_2,\mu_2)$.
Соответствующие скалярные произведения будем обозначать $\langle\cdot,\cdot\rangle_1$ и
$\langle\cdot,\cdot\rangle_2$, а нормы $||\cdot||_1$ и $||\cdot||_2$.

Например, $D_1=M_1=\{1,\ldots,L\}$ и $D_2=M_2=\{1,\ldots,K\}$, $\mu_1(\{i\}) = 1$ для всех $i=1,\ldots,L$
и $\mu_2(\{j\}) = 1$ для всех $j=1,\ldots,K$ (меры, которые каждому значению дают вес 1, называются
считающими). Тогда $L_1$ --- обычное пространство векторов длины $L$
и $L_2$ --- пространство векторов длины $K$, со стандартным евклидовым скалярным произведением.

\textit{Перед тем как читать дальше, убедитесь, что вам понятно все, о чем написано выше.}

\medskip
Теперь зададим отображение $\mathbb{G}$ из $L_2$ в $L_1$: $\mathbb{G}:\, L_2 \rightarrow L_1$.
На вход $\mathbb{G}$ поступает функция из $L_2$, а  на выходе получается функция из $L_1$.

Например, если $D_1=M_1$, $D_2=M_2$ и $L_1$, $L_2$ --- пространства векторов, то $\mathbb{G}$ --- линейное отображение,
которое вектор длины $K$ переводит в вектор длины $L$.
В качестве $\mathbb{G}$ можно взять отображение, которое задается умножением на матрицу $\mathbf{G}$ размера $L$ на $K$.
Действительно, если $\mathbf{G}\in \mathbb{R}^{L\times K}$, $X\in \mathbb{R}^{K}$,
то $\mathbf{G} X = Y \in \mathbb{R}^{L}$.

Теперь введем сопряженное отображение $\mathbb{G}^*:\, L_1 \rightarrow L_2$ такое, что
$\langle f, \mathbb{G}h \rangle_1 = \langle \mathbb{G}^*f, h \rangle_2$, где $f\in L_1$,
$h\in L_2$.

\textit{Задание 2. Покажите, что в случае пространств векторов и считающих мер, когда действие оператора
$\mathbb{G}$ --- это умножение на матрицу $\mathbf{G}$,
$\mathbb{G}^*$ --- это умножение на транспонированную матрицу $\mathbf{G}^\mathrm{T}$.}

Оператор $\mathbb{G}$, по определению, задается ядром $g(x,s)$ и действует как
$$(\mathbb{G}h)(x) = \int_{D_2} g(x,s) h(s) \mu_2(ds).$$
Требуется, чтобы
$\int_{D_1}\int_{D_2} g(x,s)^2 \mu_1(dx) \mu_2(ds)$ был конечен. В случае функций-векторов ($D_1=M_1$, $D_2=M_2$),
здесь ничего страшного нет: $g(x,s)$ --- это просто элементы матрицы, интеграл (сумма) всегда конечен,
$h$ --- вектор, действие оператора --- обычное (если мера считающая) произведение матрицы на вектор.  Такой оператор
в общем случае называется оператором
Гильберта-Шмидта. Можно думать о нем как об обобщении оператора-матрицы, переводящей умножением вектор в вектор.

\textit{Задание 3: написать, как действует оператор $\mathbb{G}$ на вектор, если меры не обязательно считающие.
Например, $\mu_1(\{i\})=w_i$, $\mu_2(\{j\})=q_j$.}

Далее рассматриваются отображения $\mathbb{G}\mathbb{G}^*: L_1\rightarrow L_1$ и
$\mathbb{G}^*\mathbb{G}: L_2\rightarrow L_2$.

Ядро оператора $\mathbb{G}\mathbb{G}^*$ выглядит следующим образом (не пугайтесь, в случае матриц это просто перемножение матриц,
но если мера не считающая, то получается похитрее):
$g_{11}(x,y) = \int_{D_2} g(x,s) g(y,s) \mu_2(ds)$.

\textit{Задание 4: выпишите в случае, когда $D_1=M_1$, $D_2=M_2$, $g(i,j)$  --- элементы матрицы $L\times K$, чему равно $\mathbb{G}\mathbb{G}^* Y$, где $Y\in \mathbb{R}^L$,
если меры не обязательно считающие: $\mu_1(\{i\})=w_1$, $\mu_2(\{j\})=q_j$.}

У каждого из этих операторов есть не более чем счетное число ненулевых (положительных)
собственных чисел и соответствующих собственных функций (определяются так же, как
для матриц). Обозначим ортонормированную систему собственных функций оператора
$\mathbb{G}\mathbb{G}^*$ как $\{\phi_n\}, \phi_n\in L_1$, ортонормированную систему собственных функций оператора
$\mathbb{G}^*\mathbb{G}$ как $\{\psi_n\}, \psi_n\in L_2$. Известно, что им соответствуют одни и те же собственные числа
$\{\lambda_n\}$, и что $\psi_n = \mathbb{G}^*\phi_n/\sqrt{\lambda_n}$, $\phi_n = \mathbb{G}\psi_n/\sqrt{\lambda_n}$.
Ортогональность означает, как обычно, равенство нулю скалярного произведения,
а нормированность --- равенство единице нормы.  Как скалярное произведение, так и норма, зависят от мер,
заданных в соответствующих подпространствах.

Обращаю внимание, что для операторов $\mathbb{G}$, заданных одним и тем же ядром
$g$, собственные числа и вектора могут различаться, если меры в пространствах заданы разные.

Основной результат про разложение Шмидта ядра оператора $\mathbb{G}$:
$$g(x,s) = \sum_n \sqrt{\lambda_n} \phi_n(x) \psi_n(s).$$

\textit{Задание 5: Покажите, что если $g(i,j)$  --- элементы матрицы  $\mathbf{G}$ размера $L\times K$, меры считающие
($\mu_1(\{i\})=1$, $\mu_2(\{j\})=1$), то разложение Шмидта ядра оператора $\mathbb{G}$ --- это в точности сингулярное разложение матрицы $\mathbf{G}$.}

\end{document}
